\SourceLeaf{11}
Es sind zugleich die Lösungen der Kongruenz
\[
1+(x+iy)^{p+1}\equiv0\quad\text{für}\quad p\equiv3\quad(\mathrm{mod.}\ 4).
\]
\par\indent Aus einem Wert $x+iy$ erhält man alle, indem man für $u$ alle Werte\\
$0,1,2,3,\, .\, .\, .,p-1$ einsetzt, und zwar mithilfe der Formel\HistoricalNote{[*)]}{[*)}{Neben dieser Herleitung stehen im Manuskript die Formeln:]
\[
1-xx\equiv yy,\qquad x=1-uy,\qquad 2u-uuy=y
\]
\[
\frac{1-uu}{1+uu}=x,\qquad\frac{2u}{1+uu}=y,
\]
\noindent [welche die Rolle des Zusatzfaktors $\dfrac{uu-1+2iu}{uu+1}$ erläutern.]}
\[
(x+iy)\left(\frac{uu-1+2iu}{uu+1}\right)\quad\text{oder}\quad\frac{(x+iy)(u+i)}{u-i}.
\]
\par\indent Für $p\equiv1\ (\mathrm{mod.}\ 4)=aa+bb$ enthält die Formel $\dfrac{b}{a}.\dfrac{u+i}{u-i}$ alle Werte von $x+iy$, wobei für $u$ nur die Werte $\dfrac{a}{b}$ und $\dfrac{b}{a}$ auszuschließen sind.
